\section{经济与社会：维持一张封建网络要付出什么}

西周的经济不能简单想象成后世那种统一市场。农业、贡纳、役使、战争缴获、赏赐和宗族控制交织在一起。王室要让远近封国持续承认自己，就必须不断发放名器、土地、职位和礼仪荣誉；这些并不只是象征，它们背后都需要粮食、劳动力、交通和军事资源。

普通人很少以自己的名字进入传世史书，但他们并未缺席：耕作、服役、筑城、运输和从军构成了上层政治能否运转的底座。青铜器铭文常记赏赐与功劳，也间接提醒我们，资源的流动首先服务于贵族网络。我们不能把后世“井田制”的整齐模型原样套回西周；更稳妥的说法是，土地、劳役和宗族权利彼此缠绕，王室与贵族都依赖对人和资源的持续动员。

\subsection{一只鼎为什么能讲经济史}

西周青铜器需要铜、锡、工匠、运输和长期的组织能力。器物铭文记录册命、赏赐、征伐和祭祀，表面是在纪念某位贵族，背后却让我们看见一套资源如何从农业和劳役中被抽取，再转化为军队、礼仪和政治记忆。

因此，金文不能只被当成“王室说了什么”。它也能让我们追问：谁有资格获得赏赐？谁能组织铸造？王室怎样把一次征伐的战功变成一个家族长期拥有的名分？这些问题把政治史和经济社会史接在一起。

\subsection{东方经营的社会成本}

\eventref{WZ-1040-DONGZHENG}{西周·三监东征}之后的东方安置，不只是把人搬到地图上。迁置、驻军、筑城、祭祀与封国重组，都需要当地社会承担新的劳役和服从关系。材料不足以让我们逐户复原这种生活，但足以提醒我们：周初秩序的形成，不只是一群贵族的会议，也是许多人被重新编入新的政治网络。

\begin{sourcenote}[经济社会史怎样避免过度推断]
传世文献偏重王室和贵族，金文偏重册命与功劳，考古则提供墓葬、聚落、作坊和器物的物质线索。三者能互相补充，但都无法直接告诉我们每个家庭的收入和感受。正文会标明哪些是材料可见的关系，哪些只是从制度运转推得的可能成本。
\end{sourcenote}

\begin{turningpoint}[制度得失：分封网络的效率与代价]
\term{它解决了什么}：在交通有限、行政工具有限的条件下，周不可能用少量王室官员直接管理广阔地区。分封让地方精英以自己的宗族和资源承担治理、守卫和扩张，降低了中心的日常成本。

\term{谁承担代价}：地方人群承担赋役和战争的直接负担；王室则要不断用赏赐、册命和婚姻维持地方精英的忠诚。中心无法免费拥有忠诚。

\term{它留下什么压力}：封国一代代积累土地、人口和军事资源后，王室与诸侯之间不再是简单的上下级。到了\eventref{WZ-0771-HAOJING}{西周·镐京陷落}，原本用来稳定边疆和东方的网络，也成为不同政治联盟各自调动资源的渠道。
\end{turningpoint}
